% ===========================================================================
\chapter{What we are going to build}

A \emph{coding agent} is a program you give a task in plain words, such as
``some of the checks fail; find out why and fix it'', and that then works
on a project by itself. It looks at the files, runs the tests, reads the
code that matters, changes it, runs the tests again, and tells you what it
did. The thinking is done by a large language model: in this tutorial,
Claude. The agent is the program around the model, and that program is
what we are going to write, in Nex, from scratch.

Here is the finished agent at work on a small project, a lending library
whose fines come out wrong. Everything the agent does is printed as it
happens: what the model says, each tool it asks for (indented), and the
first line of what each tool gave back.

\laterout{code.nex --replay sessions/fix_late_fines.json --yes}{replay}

The model never touches the project. All it can do is write text. What
makes it an agent is that some of that text is a request to use a tool:
``run \code{nex checks.nex}'', ``show me \code{loans.nex}'', ``replace
these two lines with those''. Our program carries out each request, and
sends back what happened. Then the model writes some more. The whole agent
is this loop:

\begin{shell}
conversation := [the task]
repeat
    reply := the model's answer to the conversation
    add the reply to the conversation
    if the reply asks for tools then
        run them, and add what they gave back to the conversation
    else
        the task is done
\end{shell}

That is the whole idea, and Chapter~\ref{ch:loop} writes it in about forty
lines. Most of the work, and most of this tutorial, goes into what is
around the loop: tools that are safe to hand to a model and give it useful
answers, edits that fail loudly instead of quietly doing the wrong thing,
a workspace the agent cannot reach outside of, permission before anything
changes, limits, recovery from errors, and keeping the conversation from
growing for ever.

\begin{aside}
\textbf{Where the safety lives.} The model only ever \emph{asks}. Whether a
file is read, an edit is made or a command is run is decided by our
program, every time. So the rules that keep an agent safe (which folder it
may touch, what needs your permission) are ordinary code, with ordinary
contracts, that we write and check like any other.
\end{aside}

\section{How this tutorial works}

We build the agent file by file, in the order you would write it. Every
line of every file appears in the text, in boxes. A box titled \emph{New
file} starts a file. Other boxes add to a file you already have: \emph{add at the end},
\emph{add before} a given line, or \emph{continued} (right after the box
before). Follow them and you end up with the project in
\code{examples/coding_agent} of the Nex distribution, exactly: this
tutorial is built from the same source as the code, and the build checks
that the two match.

Green boxes are \emph{try-it} scripts: small throwaway programs that show
what the code so far does. The dark boxes under them are what they print,
copied from a real run, as is every output in this tutorial.

You need no language model to follow along. Until Chapter~\ref{ch:claude},
the model is played by a script we write ourselves, which lets us test
every part of the agent with replies we choose. After that, a recorded
session plays the model's part, so the outputs in this tutorial are the
same every time it is built. To have Claude itself do the work, you need
an API key, and Chapter~\ref{ch:program} shows how to use it.

\section{What you need}

\begin{itemize}
\item Nex 0.6.5 or later, with \code{nex} on your \code{PATH}. The agent
  runs \code{nex} itself, to run the checks of the project it works on.
\item \code{bash}, for the commands the agent runs. Linux and macOS have
  it.
\item Optionally, a Claude API key from \code{console.anthropic.com}, to
  run the agent for real.
\end{itemize}

You do not need to know anything about language models. You should know
some Nex, or another language in the same family; anything new is
explained when it first comes up, and Appendix~\ref{app:api} puts the API
on one page.

\section{The plan}

\begin{center}
\begin{tabular}{@{}ll@{}}
\toprule
Chapter & What we add \\
\midrule
\ref{ch:setup}. Setting up & the project folder; text helpers \\
\ref{ch:conversation}. A conversation as data & messages, tool calls, replies \\
\ref{ch:model}. A pretend model & the \code{Model} class, and a model played by a script \\
\ref{ch:workspace}. A place to work & a workspace the agent can't reach outside of \\
\ref{ch:tools}. Giving the model hands & tools: listing and reading files \\
\ref{ch:loop}. The agent loop & the agent itself \\
\ref{ch:search}. Finding things & searching the code \\
\ref{ch:editing}. Changing things safely & editing and writing files \\
\ref{ch:commands}. Running commands & running checks, with a time limit \\
\ref{ch:permission}. Asking first & permission before anything changes \\
\ref{ch:claude}. Talking to Claude & the Messages API, over HTTP \\
\ref{ch:replay}. Record and replay & sessions you can play back \\
\ref{ch:program}. The program & \code{code.nex}, the agent on the command line \\
\ref{ch:long}. Long conversations & counting tokens; summing up and starting afresh \\
\ref{ch:checks}. Trust, but verify & executable checks for all of it \\
\ref{ch:next}. Where to go next & streaming, sub-agents, undo, \ldots \\
\bottomrule
\end{tabular}
\end{center}
